\subsection{Offsetting point-map differences on matched support}
The auxiliary whole-window increase is $\MapNet\times10^3$ Mg. At native posting, positive and negative local mass differences sum to $\MapGain\times10^3$ and $\MapLoss\times10^3$ Mg, respectively. Thus $\MapCancellation\%$ of the sum of absolute local differences cancels algebraically in the net. These are differences between published cell-level point estimates, not measured gross growth and loss or an error-adjusted ecological-change estimate. Geographic cell areas, rather than nominal equal areas, are used throughout.

Coarsening preserves the net but reduces the observable offsets (Fig.~\ref{fig:identification}a). At blocks of $64\times64$ native cells, every block difference is nonnegative and the cancellation fraction is zero, although the fine-grid negative differences remain in the underlying point maps. The result demonstrates information lost by this aggregation operation; it does not establish that every fine-grid difference is resolved by the sensor or distinguish ecological change from spatially correlated error.

\subsection{Constructive sign ambiguity and information increments}
The extremal witnesses attain the difference bounds in Table~\ref{tab:allocation} while leaving all native-cell totals and mapped fractions unchanged. In particular, the same positive whole-window difference is compatible with a formations allocation difference of $-223.97\times10^3$ Mg or $282.00\times10^3$ Mg. For the separately defined formations-plus-savanna domain the corresponding extremes are $-178.11\times10^3$ and $259.53\times10^3$ Mg. Intermediate differences, including zero, are attainable by convex interpolation. Therefore even observing an aggregate biomass increase cannot, under this allocation model, rule out a mapped-group decline. This is a test of sign identification from the available point-map information, not evidence that either extreme occurred or that the two overlapping groups can attain their separate extrema jointly.

The coherent joint identified set is a hexagon rather than the rectangle formed from the separate marginal intervals (Fig.~\ref{fig:joint}). Nevertheless, all four strict sign combinations remain feasible under a single three-component allocation at each endpoint. For example, one allocation gives formations $-119.38\times10^3$ Mg and formations plus savanna $119.38\times10^3$ Mg; another gives $103.68\times10^3$ and $-103.68\times10^3$ Mg, respectively. Thus the ambiguity is not an artifact of allowing incompatible biomass allocations to the nested groups. These are coherent hypothetical allocations, not evidence of observed group trajectories.

\begin{figure}[htbp]
\centering\includegraphics[width=.78\textwidth]{joint_identification_F11.pdf}
\caption{Sharp joint difference set for the nested mapped groups, conditional on exact point maps, nonnegative unrestricted allocation and no temporal coupling. The added-savanna constraint cuts off combinations admitted by the rectangle of marginal ranges. Four reconstructed allocations attain every strict sign quadrant while preserving all native-cell totals and component-presence geometry. Points are feasible allocation witnesses, not ecological observations or confidence regions.}\label{fig:joint}
\end{figure}

\begin{table}[htbp]
\centering\small
\caption{Same-footprint information comparison under the unrestricted allocation model. Widths are in $10^3$ Mg. Reduction compares native-cell mass observations with the one-window mass observation, retaining identical mapped domains and point maps. Neither information set identifies a group-change sign.}\label{tab:frontier}
\begin{tabular}{lrrr}
\toprule
Mapped domain & One-window width & Native-cell width & Reduction (\%)\\\midrule
Formations & 993.09 & 505.97 & 49.05\\
Formations + savanna & 993.09 & 437.64 & 55.93\\
\bottomrule
\end{tabular}
\end{table}


The one-window information set has difference bounds $[-463.92,529.16]\times10^3$ Mg for either group, because both groups and their complements occupy positive area. Restoring native-cell mass observations narrows its width by $\ReductionForm\%$ for formations and $\ReductionSav\%$ for formations plus savanna (Table~\ref{tab:frontier}). Neither narrower set identifies a sign. The nine nested observation scales show that the width has already reached the one-window value at $32\times32$-cell blocks (Fig.~\ref{fig:identification}b). This is an exact comparison of identification sets under matched support, not a held-out reporting-rule or classification-performance comparison.

\subsection{Allocation restrictions required for a positive sign}
Within the density-contrast sensitivity family, homogeneous allocation ($R=1$) gives positive model differences of $28.78\times10^3$ Mg for formations and $41.87\times10^3$ Mg for formations plus savanna. These are assumption-dependent outputs, not empirically identified forest-stock estimates. The lower difference bound reaches zero at $R=\RatioForm$ and $R=\RatioSav$, respectively (Fig.~\ref{fig:identification}c). Strict positive-sign identification within this family therefore requires a common contrast bound below the corresponding value; at the boundary zero is permitted. Neither contrast bound has been measured or validated here. The unrestricted sets remain the primary results. The sensitivity analysis identifies a specific additional information requirement rather than supplying it by assumption.

\begin{figure}[htbp]
\centering\includegraphics[width=\textwidth]{identification_frontier_F11.pdf}
\caption{Fixed-support point-map information analysis. (a) Summed positive and negative local mass differences decrease under spatial aggregation while the net is preserved. (b) Mapped-group difference-set widths increase as native-cell mass observations are replaced by block totals; cover definitions and footprint remain fixed. (c) The lower difference bound under an assumed within-cell density-contrast limit $R$; crosses mark zero-boundary values. No finite $R$ is empirically established. Native posting is not independent effective resolution. Panels describe conditional point-map arithmetic and allocation information, not calibrated ecological change, national trends or classification accuracy.}\label{fig:identification}
\end{figure}
